
\documentclass{article}
\usepackage{colortbl}
\usepackage{makecell}
\usepackage{multirow}
\usepackage{supertabular}

\begin{document}

\newcounter{utterance}

\centering \large Interaction Transcript for game `wordle', experiment `medium\_frequency\_words\_no\_clue\_no\_critic', episode 10 with gemini{-}3{-}flash{-}t1.0.
\vspace{24pt}

{ \footnotesize  \setcounter{utterance}{1}
\setlength{\tabcolsep}{0pt}
\begin{supertabular}{c@{$\;$}|p{.15\linewidth}@{}p{.15\linewidth}p{.15\linewidth}p{.15\linewidth}p{.15\linewidth}p{.15\linewidth}}
    \# & \multicolumn{2}{c}{Player} && \multicolumn{2}{c}{Game Master} \\
    \hline

    \theutterance \stepcounter{utterance}  
    & & & \multicolumn{4}{p{0.6\linewidth}}{
        \cellcolor[rgb]{0.9,0.9,0.9}{
            \makecell[{{p{\linewidth}}}]{
                \texttt{\tiny{[P1$\langle$GM]}}
                \texttt{You are a language wizard who likes to guess words by using the given rules.} \\
\\ 
\texttt{Welcome to Wordle! You have six attempts to guess the target word, a valid English word of five lowercase letters (a{-}z). Please use the tags "explanation:" and "guess:" to provide a concise explanation for each guess.} \\
\\ 
\texttt{For instance, if your guess is "apple", your response should be} \\
\texttt{explanation: this is a common five{-}letter English word, and I am starting my guess with this word.} \\
\texttt{guess: apple} \\
\\ 
\texttt{After each guess, your answer will be validated, and you will receive feedback indicating which letters are correct (green), which letters are correct but in the wrong position (yellow), and which letters are incorrect (red). This feedback can be useful in determining which letters to include or exclude in your next guess.} \\
\\ 
\texttt{For example, the feedback for "apple" might be:} \\
\texttt{guess\_feedback: a<yellow> p<yellow> p<green> l<yellow> e<red>} \\
\\ 
\texttt{The explanation should contain details about how the guess\_feedback is used to arrive at a new guess.} \\
\\ 
\texttt{Do not generate any other text.} \\
\\ 
\texttt{Let's begin with your first guess.} \\
            }
        }
    }
     \\ \\

    \theutterance \stepcounter{utterance}  
    & \multicolumn{4}{p{0.6\linewidth}}{
        \cellcolor[rgb]{0.561,0.737,0.561}{
            \makecell[{{p{\linewidth}}}]{
                \texttt{\tiny{[P1$\rangle$GM]}}
                \texttt{explanation: I am starting with "slate" as it contains common vowels and consonants to help narrow down the possibilities.} \\
\texttt{guess: slate} \\
            }
        }
    }
    & & \\ \\

    \theutterance \stepcounter{utterance}  
    & & & \multicolumn{2}{p{0.3\linewidth}}{
        \cellcolor[rgb]{0.9,0.9,0.9}{
            \makecell[{{p{\linewidth}}}]{
                \texttt{\tiny{[GM$|$GM]}}
                \texttt{attempts = 1} \\
\texttt{target\_word = tenth} \\
\texttt{guess = slate} \\
\texttt{guess\_feedback = s<red> l<red> a<red> t<green> e<yellow>} \\
            }
        }
    }
    & & \\ \\

    \theutterance \stepcounter{utterance}  
    & & & \multicolumn{4}{p{0.6\linewidth}}{
        \cellcolor[rgb]{0.9,0.9,0.9}{
            \makecell[{{p{\linewidth}}}]{
                \texttt{\tiny{[P1$\langle$GM]}}
                \texttt{guess\_feedback: s<red> l<red> a<red> t<green> e<yellow>} \\
\\ 
\texttt{Provide your response only in this format.} \\
\texttt{explanation: details} \\
\texttt{guess: word} \\
            }
        }
    }
     \\ \\

    \theutterance \stepcounter{utterance}  
    & \multicolumn{4}{p{0.6\linewidth}}{
        \cellcolor[rgb]{0.561,0.737,0.561}{
            \makecell[{{p{\linewidth}}}]{
                \texttt{\tiny{[P1$\rangle$GM]}}
                \texttt{explanation: The previous feedback shows that 't' is correct in the fourth position, and 'e' is in the word but not in the fifth position. I will try "deity" to test 'e' in the second position and introduce new letters 'd', 'i', and 'y'.} \\
\texttt{guess: deity} \\
            }
        }
    }
    & & \\ \\

    \theutterance \stepcounter{utterance}  
    & & & \multicolumn{2}{p{0.3\linewidth}}{
        \cellcolor[rgb]{0.9,0.9,0.9}{
            \makecell[{{p{\linewidth}}}]{
                \texttt{\tiny{[GM$|$GM]}}
                \texttt{attempts = 2} \\
\texttt{target\_word = tenth} \\
\texttt{guess = deity} \\
\texttt{guess\_feedback = d<red> e<green> i<red> t<green> y<red>} \\
            }
        }
    }
    & & \\ \\

    \theutterance \stepcounter{utterance}  
    & & & \multicolumn{4}{p{0.6\linewidth}}{
        \cellcolor[rgb]{0.9,0.9,0.9}{
            \makecell[{{p{\linewidth}}}]{
                \texttt{\tiny{[P1$\langle$GM]}}
                \texttt{guess\_feedback: d<red> e<green> i<red> t<green> y<red>} \\
\\ 
\texttt{Provide your response only in this format.} \\
\texttt{explanation: details} \\
\texttt{guess: word} \\
            }
        }
    }
     \\ \\

    \theutterance \stepcounter{utterance}  
    & \multicolumn{4}{p{0.6\linewidth}}{
        \cellcolor[rgb]{0.561,0.737,0.561}{
            \makecell[{{p{\linewidth}}}]{
                \texttt{\tiny{[P1$\rangle$GM]}}
                \texttt{explanation: From the feedback, 'e' is correct in the second position and 't' is correct in the fourth position. The word follows the pattern \_ e \_ t \_. "tenth" uses these confirmed positions and tests common letters 'n' and 'h' while also checking if 't' appears at the beginning of the word.} \\
\texttt{guess: tenth} \\
            }
        }
    }
    & & \\ \\

    \theutterance \stepcounter{utterance}  
    & & & \multicolumn{2}{p{0.3\linewidth}}{
        \cellcolor[rgb]{0.9,0.9,0.9}{
            \makecell[{{p{\linewidth}}}]{
                \texttt{\tiny{[GM$|$GM]}}
                \texttt{attempts = 3} \\
\texttt{target\_word = tenth} \\
\texttt{guess = tenth} \\
\texttt{guess\_feedback = t<green> e<green> n<green> t<green> h<green>} \\
            }
        }
    }
    & & \\ \\

    \theutterance \stepcounter{utterance}  
    & & & \multicolumn{2}{p{0.3\linewidth}}{
        \cellcolor[rgb]{0.9,0.9,0.9}{
            \makecell[{{p{\linewidth}}}]{
                \texttt{\tiny{[GM$|$GM]}}
                \texttt{game\_result = WIN} \\
            }
        }
    }
    & & \\ \\

\end{supertabular}
}

\end{document}
